\documentclass[12pt]{ximera}
\typeout{Start loading xmPreamble.tex}%

% Add here extra macro's that are loaded automatically by all documents of claas 'ximera' or 'xourse' in this repo

%%
%%  Example:
%%
% \newcommand{\R}{\mathbb{R}}
\usepackage{tikz}
\usetikzlibrary{positioning, arrows.meta}
\usepackage{multicol}
\usepackage{enumitem}
\usepackage{caption}
\usepackage{amsmath}
\usepackage{amsthm}
\usepackage{fontenc}

%% ---------------------------------------------------------------
%% Shared worksheet macros.
%%
%% These came from the Summer 2026 worksheet set, where each worksheet carried its
%% own copy. They have to live here instead: a xourse inlines every activity into a
%% single document, so duplicate \newcommand definitions across activities collide,
%% and the xourse gobbles \input so a per-topic macro file never loads.
%%
%% They use \providecommand, NOT \newcommand. ximera.cls loads this file automatically,
%% and every activity also carries an explicit \typeout{Start loading xmPreamble.tex}%

% Add here extra macro's that are loaded automatically by all documents of claas 'ximera' or 'xourse' in this repo

%%
%%  Example:
%%
% \newcommand{\R}{\mathbb{R}}
\usepackage{tikz}
\usetikzlibrary{positioning, arrows.meta}
\usepackage{multicol}
\usepackage{enumitem}
\usepackage{caption}
\usepackage{amsmath}
\usepackage{amsthm}
\usepackage{fontenc}

%% ---------------------------------------------------------------
%% Shared worksheet macros.
%%
%% These came from the Summer 2026 worksheet set, where each worksheet carried its
%% own copy. They have to live here instead: a xourse inlines every activity into a
%% single document, so duplicate \newcommand definitions across activities collide,
%% and the xourse gobbles \input so a per-topic macro file never loads.
%%
%% They use \providecommand, NOT \newcommand. ximera.cls loads this file automatically,
%% and every activity also carries an explicit \typeout{Start loading xmPreamble.tex}%

% Add here extra macro's that are loaded automatically by all documents of claas 'ximera' or 'xourse' in this repo

%%
%%  Example:
%%
% \newcommand{\R}{\mathbb{R}}
\usepackage{tikz}
\usetikzlibrary{positioning, arrows.meta}
\usepackage{multicol}
\usepackage{enumitem}
\usepackage{caption}
\usepackage{amsmath}
\usepackage{amsthm}
\usepackage{fontenc}

%% ---------------------------------------------------------------
%% Shared worksheet macros.
%%
%% These came from the Summer 2026 worksheet set, where each worksheet carried its
%% own copy. They have to live here instead: a xourse inlines every activity into a
%% single document, so duplicate \newcommand definitions across activities collide,
%% and the xourse gobbles \input so a per-topic macro file never loads.
%%
%% They use \providecommand, NOT \newcommand. ximera.cls loads this file automatically,
%% and every activity also carries an explicit \input{../xmPreamble}, so this file is
%% read twice. That was harmless while it held only \usepackage lines (LaTeX skips a
%% package it has already loaded) but a second \newcommand is a hard error.
%%
%%   \blank[width]        fill-in-the-blank rule (default 2.5cm)
%%   \showwork            "show and explain your work" reminder
%%   \showworkline        ... plus which schedule line was used
%%   \showworkyear        ... plus which year's schedule was used
%%   \schedhead{...}      centred bold caption above a tiered-rate table
%%   \samesquares[scale]{left}{right}   two equal-area squares, labelled
%%   \samecubes[scale]{left}{right}     two equal-volume cubes, labelled
%% ---------------------------------------------------------------

\providecommand{\blank}[1][2.5cm]{\underline{\hspace{#1}}}
\providecommand{\showwork}{\textbf{REMEMBER TO SHOW AND EXPLAIN YOUR WORK.}}
\providecommand{\showworkline}{\textbf{REMEMBER TO SHOW AND EXPLAIN YOUR WORK.
  Explanation should include how and why you know which line of the
  schedule was used to calculate this amount.}}
\providecommand{\showworkyear}{\begin{sloppypar}\textbf{REMEMBER TO SHOW AND
  EXPLAIN YOUR WORK. Explanation should include how and why you know which
  line of the year's tax schedule was used to calculate this tax
  amount.}\end{sloppypar}}
\providecommand{\schedhead}[1]{\begin{center}\textbf{#1}\end{center}}

\providecommand{\samesquares}[3][1]{%
\begin{center}
{\footnotesize These squares are the \textbf{same size}.}

\vspace{1.5mm}
\begin{tikzpicture}[scale=#1,every node/.style={scale=#1}]
  \draw (0,0) rectangle (2.4,2.4);
  \node[anchor=north west] at (0.15,2.25) {Area:};
  \node[anchor=east]  at (-0.15,1.2) {#2};
  \node[anchor=north] at (1.2,-0.15) {#2};
  \begin{scope}[xshift=4.6cm]
    \draw (0,0) rectangle (2.4,2.4);
    \node[anchor=north west] at (0.15,2.25) {Area:};
    \node[anchor=east]  at (-0.15,1.2) {#3};
    \node[anchor=north] at (1.2,-0.15) {#3};
  \end{scope}
\end{tikzpicture}
\end{center}}

\providecommand{\samecubes}[3][1]{%
\begin{center}
{\footnotesize These cubes are the \textbf{same size}.}

\vspace{1.5mm}
\begin{tikzpicture}[scale=#1,every node/.style={scale=#1}]
  \foreach \dx in {0,5.2} {
    \begin{scope}[xshift=\dx cm]
      \draw (0,0) rectangle (2.0,2.0);
      \draw (0,2.0) -- (0.65,2.6) -- (2.65,2.6) -- (2.0,2.0);
      \fill[black!12] (2.0,0) -- (2.65,0.6) -- (2.65,2.6) -- (2.0,2.0) -- cycle;
      \draw (2.0,0) -- (2.65,0.6) -- (2.65,2.6) -- (2.0,2.0);
      \node[anchor=north west] at (0.1,1.9) {Volume:};
    \end{scope}
  }
  \node[anchor=east]  at (-0.15,1.0) {#2};
  \node[anchor=north] at (1.0,-0.15) {#2};
  \node[anchor=west]  at (2.25,0.4) {#2};
  \node[anchor=east]  at (5.05,1.0) {#3};
  \node[anchor=north] at (6.2,-0.15) {#3};
  \node[anchor=west]  at (7.45,0.4) {#3};
\end{tikzpicture}
\end{center}}

%% NOTE: do NOT add \usepackage to individual activity .tex files.
%% When a xourse inlines an activity it gobbles \documentclass and \input, but a bare
%% \usepackage still executes -- after \begin{document} -- and errors out.
%% All shared packages and macros belong here.
%%
%% ximera.cls already loads: amsmath, amssymb, amsthm, cancel, enumitem, graphicx,
%% hyperref, listings, pgfplots, tikz, url, xcolor. No need to re-load those.
, so this file is
%% read twice. That was harmless while it held only \usepackage lines (LaTeX skips a
%% package it has already loaded) but a second \newcommand is a hard error.
%%
%%   \blank[width]        fill-in-the-blank rule (default 2.5cm)
%%   \showwork            "show and explain your work" reminder
%%   \showworkline        ... plus which schedule line was used
%%   \showworkyear        ... plus which year's schedule was used
%%   \schedhead{...}      centred bold caption above a tiered-rate table
%%   \samesquares[scale]{left}{right}   two equal-area squares, labelled
%%   \samecubes[scale]{left}{right}     two equal-volume cubes, labelled
%% ---------------------------------------------------------------

\providecommand{\blank}[1][2.5cm]{\underline{\hspace{#1}}}
\providecommand{\showwork}{\textbf{REMEMBER TO SHOW AND EXPLAIN YOUR WORK.}}
\providecommand{\showworkline}{\textbf{REMEMBER TO SHOW AND EXPLAIN YOUR WORK.
  Explanation should include how and why you know which line of the
  schedule was used to calculate this amount.}}
\providecommand{\showworkyear}{\begin{sloppypar}\textbf{REMEMBER TO SHOW AND
  EXPLAIN YOUR WORK. Explanation should include how and why you know which
  line of the year's tax schedule was used to calculate this tax
  amount.}\end{sloppypar}}
\providecommand{\schedhead}[1]{\begin{center}\textbf{#1}\end{center}}

\providecommand{\samesquares}[3][1]{%
\begin{center}
{\footnotesize These squares are the \textbf{same size}.}

\vspace{1.5mm}
\begin{tikzpicture}[scale=#1,every node/.style={scale=#1}]
  \draw (0,0) rectangle (2.4,2.4);
  \node[anchor=north west] at (0.15,2.25) {Area:};
  \node[anchor=east]  at (-0.15,1.2) {#2};
  \node[anchor=north] at (1.2,-0.15) {#2};
  \begin{scope}[xshift=4.6cm]
    \draw (0,0) rectangle (2.4,2.4);
    \node[anchor=north west] at (0.15,2.25) {Area:};
    \node[anchor=east]  at (-0.15,1.2) {#3};
    \node[anchor=north] at (1.2,-0.15) {#3};
  \end{scope}
\end{tikzpicture}
\end{center}}

\providecommand{\samecubes}[3][1]{%
\begin{center}
{\footnotesize These cubes are the \textbf{same size}.}

\vspace{1.5mm}
\begin{tikzpicture}[scale=#1,every node/.style={scale=#1}]
  \foreach \dx in {0,5.2} {
    \begin{scope}[xshift=\dx cm]
      \draw (0,0) rectangle (2.0,2.0);
      \draw (0,2.0) -- (0.65,2.6) -- (2.65,2.6) -- (2.0,2.0);
      \fill[black!12] (2.0,0) -- (2.65,0.6) -- (2.65,2.6) -- (2.0,2.0) -- cycle;
      \draw (2.0,0) -- (2.65,0.6) -- (2.65,2.6) -- (2.0,2.0);
      \node[anchor=north west] at (0.1,1.9) {Volume:};
    \end{scope}
  }
  \node[anchor=east]  at (-0.15,1.0) {#2};
  \node[anchor=north] at (1.0,-0.15) {#2};
  \node[anchor=west]  at (2.25,0.4) {#2};
  \node[anchor=east]  at (5.05,1.0) {#3};
  \node[anchor=north] at (6.2,-0.15) {#3};
  \node[anchor=west]  at (7.45,0.4) {#3};
\end{tikzpicture}
\end{center}}

%% NOTE: do NOT add \usepackage to individual activity .tex files.
%% When a xourse inlines an activity it gobbles \documentclass and \input, but a bare
%% \usepackage still executes -- after \begin{document} -- and errors out.
%% All shared packages and macros belong here.
%%
%% ximera.cls already loads: amsmath, amssymb, amsthm, cancel, enumitem, graphicx,
%% hyperref, listings, pgfplots, tikz, url, xcolor. No need to re-load those.
, so this file is
%% read twice. That was harmless while it held only \usepackage lines (LaTeX skips a
%% package it has already loaded) but a second \newcommand is a hard error.
%%
%%   \blank[width]        fill-in-the-blank rule (default 2.5cm)
%%   \showwork            "show and explain your work" reminder
%%   \showworkline        ... plus which schedule line was used
%%   \showworkyear        ... plus which year's schedule was used
%%   \schedhead{...}      centred bold caption above a tiered-rate table
%%   \samesquares[scale]{left}{right}   two equal-area squares, labelled
%%   \samecubes[scale]{left}{right}     two equal-volume cubes, labelled
%% ---------------------------------------------------------------

\providecommand{\blank}[1][2.5cm]{\underline{\hspace{#1}}}
\providecommand{\showwork}{\textbf{REMEMBER TO SHOW AND EXPLAIN YOUR WORK.}}
\providecommand{\showworkline}{\textbf{REMEMBER TO SHOW AND EXPLAIN YOUR WORK.
  Explanation should include how and why you know which line of the
  schedule was used to calculate this amount.}}
\providecommand{\showworkyear}{\begin{sloppypar}\textbf{REMEMBER TO SHOW AND
  EXPLAIN YOUR WORK. Explanation should include how and why you know which
  line of the year's tax schedule was used to calculate this tax
  amount.}\end{sloppypar}}
\providecommand{\schedhead}[1]{\begin{center}\textbf{#1}\end{center}}

\providecommand{\samesquares}[3][1]{%
\begin{center}
{\footnotesize These squares are the \textbf{same size}.}

\vspace{1.5mm}
\begin{tikzpicture}[scale=#1,every node/.style={scale=#1}]
  \draw (0,0) rectangle (2.4,2.4);
  \node[anchor=north west] at (0.15,2.25) {Area:};
  \node[anchor=east]  at (-0.15,1.2) {#2};
  \node[anchor=north] at (1.2,-0.15) {#2};
  \begin{scope}[xshift=4.6cm]
    \draw (0,0) rectangle (2.4,2.4);
    \node[anchor=north west] at (0.15,2.25) {Area:};
    \node[anchor=east]  at (-0.15,1.2) {#3};
    \node[anchor=north] at (1.2,-0.15) {#3};
  \end{scope}
\end{tikzpicture}
\end{center}}

\providecommand{\samecubes}[3][1]{%
\begin{center}
{\footnotesize These cubes are the \textbf{same size}.}

\vspace{1.5mm}
\begin{tikzpicture}[scale=#1,every node/.style={scale=#1}]
  \foreach \dx in {0,5.2} {
    \begin{scope}[xshift=\dx cm]
      \draw (0,0) rectangle (2.0,2.0);
      \draw (0,2.0) -- (0.65,2.6) -- (2.65,2.6) -- (2.0,2.0);
      \fill[black!12] (2.0,0) -- (2.65,0.6) -- (2.65,2.6) -- (2.0,2.0) -- cycle;
      \draw (2.0,0) -- (2.65,0.6) -- (2.65,2.6) -- (2.0,2.0);
      \node[anchor=north west] at (0.1,1.9) {Volume:};
    \end{scope}
  }
  \node[anchor=east]  at (-0.15,1.0) {#2};
  \node[anchor=north] at (1.0,-0.15) {#2};
  \node[anchor=west]  at (2.25,0.4) {#2};
  \node[anchor=east]  at (5.05,1.0) {#3};
  \node[anchor=north] at (6.2,-0.15) {#3};
  \node[anchor=west]  at (7.45,0.4) {#3};
\end{tikzpicture}
\end{center}}

%% NOTE: do NOT add \usepackage to individual activity .tex files.
%% When a xourse inlines an activity it gobbles \documentclass and \input, but a bare
%% \usepackage still executes -- after \begin{document} -- and errors out.
%% All shared packages and macros belong here.
%%
%% ximera.cls already loads: amsmath, amssymb, amsthm, cancel, enumitem, graphicx,
%% hyperref, listings, pgfplots, tikz, url, xcolor. No need to re-load those.

\title{Exponentials: Growth}
\begin{document}
\begin{abstract}
Building an exponential function from a repeated percentage increase, identifying what
$a$ and $R$ mean in $y = a \cdot R^x$, contrasting constant \emph{absolute} change
(linear) with constant \emph{relative} change (exponential), and using Guess \& Check ---
and then $n$th roots --- to answer doubling-time questions.
\end{abstract}
\maketitle

\section*{Where Exponentials Come From}

What is an exponential function? What does it mean for something to be ``exponential''?
To get an idea, let's revisit a problem from our Percentages activity.

An athlete signs a contract saying that they will earn \$8.3 million with an increase of
$4.8\%$ each year of the 5 year contract.

\begin{problem}
Fill in the athlete's salary, in millions of dollars, after each year.

After year 1: \$$\answer[tolerance=0.1]{8.7}$ million

After year 2: \$$\answer[tolerance=0.1]{9.1}$ million

After year 3: \$$\answer[tolerance=0.1]{9.55}$ million

After year 4: \$$\answer[tolerance=0.1]{10.0}$ million

\begin{solution}
Each year the salary is multiplied by $1 + 0.048 = 1.048$. Working from the previous
year's salary each time:
\begin{align*}
\text{Year 1:} \quad & \$8.3 \cdot 0.048 + \$8.3 = \$8.3 \cdot 1.048 = \$8.70 \\
\text{Year 2:} \quad & \$8.70 \cdot 1.048 = [\$8.3 \cdot 1.048] \cdot 1.048 = \$9.12 \\
\text{Year 3:} \quad & \$9.12 \cdot 1.048 = \$8.3 \cdot (1.048)^2 \cdot 1.048 = \$9.55 \\
\text{Year 4:} \quad & \$9.55 \cdot 1.048 = \$8.3 \cdot (1.048)^3 \cdot 1.048 = \$10.01
\end{align*}

But look at the right-hand side of each line. Each one collapses to the starting salary
times a power of $1.048$:
\[
\$8.3 \cdot (1.048)^0, \quad \$8.3 \cdot (1.048)^1, \quad \$8.3 \cdot (1.048)^2,
\quad \$8.3 \cdot (1.048)^3, \quad \$8.3 \cdot (1.048)^4,
\]
and the exponent is always the number of years completed. That is the pattern the whole
activity rests on.
\end{solution}
\end{problem}

\begin{problem}
Suppose this athlete renewed this contract indefinitely, so the $4.8\%$ yearly increase
didn't stop after 5 years. How much would they be earning at the end of the 17th year?

\$$\answer[tolerance=0.3]{18.4}$ million

\begin{solution}
The power always equals how many years of the contract have been completed, so put 17 in
as the exponent:
\[
\$8.3 \text{ mill} \cdot (1.048)^{17} = \$18.4 \text{ mill}.
\]
\end{solution}
\end{problem}

\begin{problem}
How much money would this athlete be earning after the $x$th year? Write a function.

\begin{freeResponse}
\end{freeResponse}

\begin{solution}
Since the exponent is the number of years completed, put $x$ in for the exponent:
\[
y = \$8.3 \text{ mill} \cdot (1.048)^x.
\]
\end{solution}
\end{problem}

\section*{Reading the Pieces of $y = a \cdot R^x$}

We have developed an exponential function describing this athlete's salary. Let's
identify what each piece represents in $y = 8.3(1.048)^x$.

\begin{problem}
For the function $y = 8.3(1.048)^x$, identify the units and meaning of each piece.

\begin{freeResponse}
\end{freeResponse}

\begin{solution}
\textbf{8.3} --- units of millions of dollars. Represents the starting salary of the
contract.

\textbf{1.048} --- no units, technically, as it comes from a percentage. The growth rate
is a $4.8\%$ increase; the 1.048 represents that $4.8\%$ increase in salary each year.

\textbf{$x$} --- units of years. Represents the amount of time that has passed, that is,
which year of the contract we are after.

\textbf{$y$} --- units of millions of dollars. Represents the salary of the contract
after year $x$.
\end{solution}
\end{problem}

\section*{Comparing Exponential and Linear Functions}

Let's change the problem slightly and compare the exponential situation with a linear
one. An athlete signs a contract saying they will earn \$8.3 million with an increase of
\textbf{\$500,000} each year of the 5 year contract.

\begin{problem}
Write a linear function to represent how much money the athlete will be making after year
$x$.

\begin{freeResponse}
\end{freeResponse}

\begin{solution}
A linear function has the form $y = mx + b$, so we need an $m$ and a $b$.

The $b$-value is \$8.3 million, as that is what the athlete is paid initially. The
\$500,000 per year is the slope, as that is how the salary changes each year.

One catch: 8.3 is in millions of dollars but 500,000 is not, so we must put them in the
same units --- either write \$8.3 mill as \$8,300,000, or write \$500,000 as \$0.5 mill:
\[
y = \frac{\$0.5 \text{ mill}}{\text{year}} \cdot x + \$8.3 \text{ mill}
\qquad \text{or} \qquad
y = \frac{\$500{,}000}{\text{year}} \cdot x + \$8{,}300{,}000.
\]
\end{solution}
\end{problem}

\begin{problem}
What is the slope of your function, what are its units, and what does it mean in this
context? Same for the $y$-intercept.

\begin{freeResponse}
\end{freeResponse}

\begin{solution}
\textbf{Slope} $= \dfrac{\$0.5 \text{ mill}}{\text{year}} = \dfrac{\$500{,}000}{\text{year}}$,
meaning the athlete will earn an additional \$0.5 million per year of the contract.

\textbf{$y$-intercept} $=$ \$8.3 mill $=$ \$8,300,000. This means that when $x = 0$ ---
at the beginning of the contract, after 0 years --- the athlete's salary is \$8.3 million.
\end{solution}
\end{problem}

\begin{problem}
Look back at the exponential function $y = 8.3(1.048)^x$. Which value in it is most like
the \emph{$y$-intercept} of the linear function, and which is most like the \emph{slope}?
Why?

\begin{freeResponse}
\end{freeResponse}

\begin{solution}
The \textbf{8.3} in the exponential function is most similar to the $y$-intercept. Not
only are they the exact same value, they both represent an initial amount.

The \textbf{1.048} is most like the slope. These values are not the same, but they are
both describing the change happening each year to the athlete's salary.

The crucial distinction: \emph{linear functions have a constant absolute change, whereas
exponential functions have a constant relative change.} The slope is a constant absolute
change --- \$500,000 more every year, no matter how big the salary already is. The
$R$-value is a constant relative change --- $4.8\%$ more every year, which is a larger
and larger dollar amount as the salary grows.
\end{solution}
\end{problem}

\section*{Building a Function Directly}

Now that we know what each piece of $y = a \cdot R^x$ means, we can write exponential
functions for new scenarios without going through the step-by-step table.

Some CDs (a banking investment option) offer rates giving about $3.10\%$ yield per year,
as long as a minimum amount such as \$5,000 is invested. A CD is held for a certain number
of years, and when it ``matures'' you usually have the option to renew.

\begin{problem}
Suppose you invest the minimum \$5,000 to get this $3.10\%$ yield. Write a function
describing how much the investment will be worth after $x$ years, and identify what each
value represents.

\begin{freeResponse}
\end{freeResponse}

\begin{solution}
\[
y = \$5{,}000 \cdot (1 + 0.031)^x = \$5{,}000(1.031)^x
\]

\$5,000 --- the initial investment, the beginning amount.

1.031 --- the growth rate, describing how the investment changes. It shows a $3.1\%$
increase for each year.

$x$ --- in years, the number of years completed in the investment.

$y$ --- in dollars, the value of the investment after $x$ years.
\end{solution}
\end{problem}

\begin{problem}
Suppose this is a 7 year CD. How much will your investment be worth at the end of the CD?

\$$\answer[tolerance=5]{6191.28}$

\begin{solution}
Since $x$ is in years, plug in 7:
\[
y = \$5{,}000 \cdot (1.031)^7 = \$6{,}191.28.
\]
\end{solution}
\end{problem}

\begin{problem}
Suppose you keep renewing this CD every time it matures. About how many years will it
take to double the initial investment? Use Guess \& Check.

$\answer{23}$ years

\begin{freeResponse}
\end{freeResponse}

\begin{solution}
We want our function to output \$10,000, because that is double the initial \$5,000
investment.

\textbf{Step 1: find the closest whole number giving less than we want.} Try whole
numbers for $x$ until you get a value below \$10,000 where the next whole number goes
above it. Here $x = 22$:
\[
\$5{,}000 \cdot (1.031)^{22} = \$9{,}787.25.
\]

\textbf{Step 2: find the closest whole number giving more than we want.} $x = 23$:
\[
\$5{,}000 \cdot (1.031)^{23} = \$10{,}090.65.
\]
So the exact answer is somewhere between $x = 22$ and $x = 23$.

\textbf{Step 3: which is closer?}
\[
\$10{,}000 - \$9{,}787.25 = \$212.75
\qquad
\$10{,}090.65 - \$10{,}000 = \$90.65
\]
\$10,090.65 is closer, so we use \textbf{23 years} as the estimate.

Even when it seems obvious which value is closer, showing this subtraction is a good
habit, and you should do it every time. In future problems asking for Guess \& Check
these steps will not be listed out for you --- they define how much work is needed for
full credit.
\end{solution}
\end{problem}

\section*{Reading Someone Else's Model}

Using data from the CDC about ebola cases in Guinea, an exponential model can be created
representing the number of cases $x$ days after May 25th, 2014, through the remainder of
2014:
\[
y = 121.2(1.0117)^x.
\]

\begin{problem}
What do $x$, $y$, and $121.2$ represent in this context?

\begin{freeResponse}
\end{freeResponse}

\begin{solution}
$x$ represents the number of days since May 25th, 2014.

$y$ represents the number of ebola cases $x$ days after May 25th, 2014.

$121.2$ is the initial amount, meaning there were about 121 cases of ebola around May
25th, 2014.
\end{solution}
\end{problem}

\begin{problem}
What does $1.0117$ tell us about this context? Give the percentage.

$\answer[tolerance=0.05]{1.17}\%$ increase per day

\begin{freeResponse}
\end{freeResponse}

\begin{solution}
This value tells us how the quantity is changing. In the previous problems we had to
build this value ourselves: for the athlete it was $1.048$ from a $4.8\%$ increase, and
for the investment it was $1.031$ from a $3.10\%$ increase.

Here we read it in the other direction: $1.0117 = 1 + 0.0117$, indicating a $1.17\%$
increase in ebola cases each day.
\end{solution}
\end{problem}

\begin{problem}
Use Guess \& Check to estimate when the number of ebola cases doubles. Show the same
steps as before.

About $\answer{60}$ days

\begin{freeResponse}
\end{freeResponse}

\begin{solution}
Double of $121.2$ is $242.4$, so we want an output around 242.4 cases.

Closest whole $x$ below: $x = 59$,
\[
121.2 \cdot (1.0117)^{59} = 240.7 \text{ cases}.
\]
Closest whole $x$ above: $x = 60$,
\[
121.2 \cdot (1.0117)^{60} = 243.6 \text{ cases}.
\]
So the exact answer is between 59 and 60 days. Which is closer?
\[
242.4 - 240.7 = 1.7 \qquad 243.6 - 242.4 = 1.2
\]
243.6 is slightly closer, so use \textbf{60 days} as the estimate.

Worth pausing on what this means: a $1.17\%$ daily increase sounds tiny, but it doubles
the caseload in two months.
\end{solution}
\end{problem}

\section*{Solving for the Rate}

\begin{problem}
Suppose there is a new virus that reportedly doubles infection cases in about 20 days.
What would be this virus's infection rate, as a percent?

About $\answer[tolerance=0.3]{3.5}\%$

\begin{freeResponse}
\end{freeResponse}

\begin{solution}
We don't know the rate, so that is the missing piece --- call it $R$. The number of days
$x$ is 20. We don't know the initial number of cases, but it will not matter: call it $a$
and set the function equal to $2a$, double whatever the initial amount is.
\[
2a = a \cdot R^{20}
\]
Divide both sides by $a$:
\[
2 = R^{20}
\]
Take the 20th root of both sides:
\[
R = 2^{1/20} = 1.0353\ldots
\]
$R = 1.035\ldots$ means an increase of about $3.5\%$, so $3.5\%$ is the infection rate.

Why the initial number wasn't needed: doubling is a \emph{relative} statement. It asks
how long until the quantity is twice whatever it started at, and the starting amount
divides out of both sides. This is a general feature of exponential growth --- the
doubling time depends only on the rate, never on where you start.
\end{solution}
\end{problem}

\begin{problem}
Return to the athlete's salary problem: \$8.3 million with a $4.8\%$ increase each year.
Determine through Guess \& Check an estimate for how long it would take for the salary to
double.

About $\answer{15}$ years

\begin{freeResponse}
\end{freeResponse}

\begin{solution}
We want the salary to be double \$8.3 mill, so we want an output of \$16.6 mill from
$y = 8.3(1.048)^x$.

Closest whole below: $x = 14$,
\[
8.3(1.048)^{14} = \$16.0 \text{ mill}.
\]
Closest whole above: $x = 15$,
\[
8.3(1.048)^{15} = \$16.8 \text{ mill}.
\]
So the answer is between 14 and 15 years. Which is closer?
\[
\$16.6 - \$16.0 = \$0.6 \text{ mill}
\qquad
\$16.8 - \$16.6 = \$0.2 \text{ mill}
\]
\$16.8 mill is closer, so the salary would double after about \textbf{15 years}.
\end{solution}
\end{problem}

\begin{problem}
Now assume the athlete negotiates for a better rate: they want their initial salary to
double in 10 years. What rate should they push for in the contract?

About $\answer[tolerance=0.3]{7.2}\%$ per year

\begin{freeResponse}
\end{freeResponse}

\begin{solution}
We don't know the rate, so put $R$ in for it. We want the time to be 10 years, so plug
that in for $x$, and set it equal to \$16.6 mill, double the initial \$8.3 mill:
\[
8.3 \cdot R^{10} = 16.6
\]
Divide both sides by 8.3:
\[
R^{10} = 2
\]
Take the 10th root of both sides:
\[
R = 2^{1/10} = 1.0718\ldots
\]
$R = 1.072\ldots$ means an increase of about $7.2\%$, so the athlete should ask for about
a $7.2\%$ increase per year.

Notice this is only about one and a half times the original $4.8\%$, yet it cuts the
doubling time from 15 years to 10. Small changes in an exponential rate have
disproportionate effects on doubling time.
\end{solution}
\end{problem}

\end{document}
